\documentclass[11pt,a4paper]{article}

%============================================================
% COURS COMPLET — ONDES ET SIGNAUX
% Première générale — Spécialité Physique-Chimie
%============================================================

\usepackage[utf8]{inputenc}
\usepackage[T1]{fontenc}
\usepackage[french]{babel}
\usepackage{lmodern}
\usepackage[a4paper,margin=2cm]{geometry}
\usepackage{amsmath,amssymb,amsthm}
\usepackage{mathtools}
\usepackage{siunitx}
\usepackage{physics}
\usepackage{fancyhdr}
\usepackage{lastpage}
\usepackage{enumitem}
\usepackage{array,tabularx,booktabs}
\usepackage{xcolor}
\usepackage{tikz}
\usepackage{pgfplots}
\usepackage{tcolorbox}
\usepackage{hyperref}
\usepackage{multicol}

\pgfplotsset{compat=1.18}
\usetikzlibrary{arrows.meta,calc,decorations.pathmorphing,patterns,positioning}
\tcbuselibrary{theorems,skins,breakable}

%============================================================
% COULEURS ET STYLE
%============================================================
\definecolor{bleuprepa}{RGB}{20,55,110}
\definecolor{vertprepa}{RGB}{20,105,70}
\definecolor{rougeprepa}{RGB}{160,35,35}
\definecolor{orangeprepa}{RGB}{180,100,20}
\definecolor{violetprepa}{RGB}{90,45,130}
\definecolor{grisclair}{RGB}{245,247,250}
\definecolor{grisfonce}{RGB}{70,70,70}

\hypersetup{
  colorlinks=true,
  linkcolor=bleuprepa,
  urlcolor=bleuprepa,
  citecolor=bleuprepa
}

\sisetup{
  locale=FR,
  per-mode=symbol,
  detect-all=true
}

\setlength{\parindent}{0pt}
\setlength{\headheight}{14pt}
\setlength{\parskip}{0.45em}
\renewcommand{\arraystretch}{1.25}

\pagestyle{fancy}
\fancyhf{}
\fancyhead[L]{Première spécialité physique-chimie}
\fancyhead[C]{Ondes et signaux}
\fancyhead[R]{\thepage/\pageref{LastPage}}
\fancyfoot[C]{Cours professeur — document pédagogique complet}

%============================================================
% ENVIRONNEMENTS
%============================================================
\tcbset{
  enhanced,
  breakable,
  sharp corners=downhill,
  boxrule=0.8pt,
  colback=white,
  fonttitle=\bfseries,
  before skip=8pt,
  after skip=8pt
}

\newtcbtheorem[number within=section]{definitionbox}{Définition}%
{colframe=bleuprepa,colback=blue!3!white}{def}

\newtcbtheorem[number within=section]{propositionbox}{Propriété}%
{colframe=vertprepa,colback=green!3!white}{prop}

\newtcbtheorem[number within=section]{theorembox}{Théorème}%
{colframe=violetprepa,colback=violet!3!white}{thm}

\newtcbtheorem[number within=section]{methodbox}{Méthode}%
{colframe=orangeprepa,colback=orange!4!white}{meth}

\newtcbtheorem[number within=section]{examplebox}{Exemple corrigé}%
{colframe=grisfonce,colback=grisclair}{ex}

\newtcbtheorem[number within=section]{remarquebox}{Remarque pédagogique}%
{colframe=bleuprepa!70!black,colback=blue!2!white}{rem}

\newtcbtheorem[number within=section]{warningbox}{Erreur classique}%
{colframe=rougeprepa,colback=red!3!white}{err}

\newtcbtheorem[number within=section]{retenirbox}{À retenir}%
{colframe=vertprepa!80!black,colback=green!5!white}{ret}

\newcounter{exercice}
\newenvironment{exercice}[1][]{%
  \refstepcounter{exercice}%
  \begin{tcolorbox}[colframe=bleuprepa,colback=white,title={Exercice \theexercice\; #1}]%
}{\end{tcolorbox}}

\newcounter{correction}
\newenvironment{correction}[1][]{%
  \refstepcounter{correction}%
  \begin{tcolorbox}[colframe=vertprepa,colback=green!2!white,title={Correction de l'exercice \thecorrection\; #1}]%
}{\end{tcolorbox}}

\renewcommand{\dd}{\,\mathrm{d}}
\newcommand{\e}{\mathrm{e}}
\newcommand{\ii}{\mathrm{i}}
\newcommand{\R}{\mathbb{R}}
\newcommand{\vect}[1]{\overrightarrow{#1}}
\newcommand{\norme}[1]{\left\lVert #1\right\rVert}

%============================================================
% DOCUMENT
%============================================================
\begin{document}

\begin{titlepage}
\centering
\vspace*{1cm}
{\Large Première générale — Spécialité Physique-Chimie\par}
\vspace{1cm}
{\Huge\bfseries Ondes et signaux\par}
\vspace{0.4cm}
{\LARGE 1. Ondes mécaniques\par}
{\LARGE 2. Lumière : images et couleurs, modèles ondulatoire et particulaire\par}
\vspace{1cm}
\begin{tcolorbox}[colframe=bleuprepa,colback=blue!3!white,width=0.92\textwidth]
\textbf{Objectif du polycopié.} Ce document propose un cours complet, progressif et rigoureux sur le thème \emph{Ondes et signaux} du programme de première spécialité physique-chimie. Il est conçu pour des élèves de première de bon niveau, mais aussi comme support professeur : définitions, propriétés, démonstrations accessibles, méthodes, exemples corrigés, exercices progressifs, corrections détaillées, devoir bilan et fiche de synthèse.
\end{tcolorbox}
\vfill
{\large Version professeur — prêt à compiler en \LaTeX\par}
\end{titlepage}

\tableofcontents
\newpage

%============================================================
\section*{Préambule : objectifs et prérequis}
\addcontentsline{toc}{section}{Préambule : objectifs et prérequis}

\begin{retenirbox}{Prérequis de seconde et de mathématiques}{pre}
Avant d'aborder ce chapitre, on doit maîtriser :
\begin{itemize}[label=$\triangleright$]
  \item la notion de signal sonore, de propagation et de vitesse de propagation ;
  \item les grandeurs période $T$ et fréquence $f$ d'un signal périodique ;
  \item la lecture d'un graphique temporel ou spatial ;
  \item les fonctions sinus et cosinus, leurs graphes et leur périodicité ;
  \item la proportionnalité, les unités, les conversions et l'écriture scientifique ;
  \item le théorème de Thalès et les grandeurs algébriques en optique.
\end{itemize}
\end{retenirbox}

\begin{remarquebox}{Philosophie du chapitre}{philo}
Le thème \emph{Ondes et signaux} articule deux idées majeures. D'abord, une perturbation peut transporter de l'information et de l'énergie sans transport global de matière : c'est le modèle ondulatoire. Ensuite, la lumière peut être décrite par plusieurs modèles complémentaires : le modèle géométrique pour former des images, le modèle ondulatoire pour relier fréquence et longueur d'onde, et le modèle particulaire pour interpréter absorption et émission de photons.
\end{remarquebox}

\newpage

%============================================================
\part{Ondes mécaniques}
%============================================================

\section{Onde mécanique progressive}

\subsection{Perturbation mécanique et milieu matériel}

\begin{definitionbox}{Perturbation mécanique}{perturbation}
Une \textbf{perturbation mécanique} est une modification locale et temporaire de l'état d'un milieu matériel. Elle peut concerner, par exemple, la position de points d'une corde, la pression de l'air, la hauteur de la surface de l'eau, ou la compression d'un ressort.
\end{definitionbox}

\begin{definitionbox}{Onde mécanique progressive}{ondeprogressive}
Une \textbf{onde mécanique progressive} est le phénomène de propagation d'une perturbation mécanique dans un milieu matériel, sans transport global de matière, mais avec transport d'énergie et d'information.
\end{definitionbox}

\begin{examplebox}{Exemples fondamentaux}{exemples-ondes}
\begin{tabularx}{\textwidth}{>{\bfseries}lX}
Onde sur une corde & une déformation de la corde se propage le long de celle-ci ; chaque point de la corde monte puis redescend localement.\\
Onde sonore & une succession de compressions et de dilatations de l'air se propage ; les molécules d'air oscillent autour d'une position moyenne.\\
Houle & la surface de l'eau est déformée ; l'énergie se propage, mais une bouée effectue surtout un mouvement local.\\
Onde sismique & une perturbation mécanique se propage dans les roches de la Terre.\\
\end{tabularx}
\end{examplebox}

\begin{warningbox}{Onde et déplacement de matière}{matiere}
Une onde ne correspond pas à un transport global de matière. Dans une corde, chaque point de la corde oscille localement, mais la corde entière ne part pas avec l'onde. Dans l'air, le son ne transporte pas les molécules depuis la source jusqu'à l'oreille : il transporte une variation de pression.
\end{warningbox}

\subsection{Modèle qualitatif de propagation}

On peut comprendre qualitativement la propagation d'une onde mécanique par l'existence d'interactions entre éléments voisins du milieu. Quand une partie du milieu est perturbée, elle agit sur ses voisines, qui deviennent perturbées à leur tour. La perturbation se transmet donc de proche en proche.

\begin{propositionbox}{Condition d'existence d'une onde mécanique}{condition-onde-meca}
Une onde mécanique nécessite :
\begin{enumerate}[label=\arabic*)]
  \item un \textbf{milieu matériel} susceptible d'être perturbé ;
  \item des \textbf{interactions} entre parties voisines du milieu ;
  \item une \textbf{source} qui crée la perturbation initiale.
\end{enumerate}
Ainsi, le son ne se propage pas dans le vide, tandis que la lumière le peut.
\end{propositionbox}

\begin{remarquebox}{Pourquoi parle-t-on de modèle ?}{modele}
Dire qu'une onde se propage de proche en proche est un modèle qualitatif. Il ne décrit pas tous les détails microscopiques du milieu, mais il explique efficacement pourquoi la perturbation avance et pourquoi la célérité dépend du milieu.
\end{remarquebox}

\subsection{Ondes transversales et longitudinales}

\begin{definitionbox}{Onde transversale et onde longitudinale}{trans-long}
On compare la direction de propagation de l'onde et la direction de la perturbation.
\begin{itemize}
  \item Une onde est \textbf{transversale} lorsque la perturbation est perpendiculaire à la direction de propagation.
  \item Une onde est \textbf{longitudinale} lorsque la perturbation est parallèle à la direction de propagation.
\end{itemize}
\end{definitionbox}

\begin{center}
\begin{tikzpicture}[scale=1]
\draw[->,thick] (0,0) -- (5.5,0) node[right]{propagation};
\draw[thick,domain=0:5,samples=100] plot (\x,{0.5*sin(360*\x/2)});
\draw[->,rougeprepa,thick] (1.25,0) -- (1.25,0.8) node[above]{perturbation};
\node at (2.8,-1) {Onde transversale sur une corde};
\begin{scope}[xshift=8cm]
\draw[->,thick] (0,0) -- (5.5,0) node[right]{propagation};
\foreach \x in {0,0.35,...,5.2}{\draw[bleuprepa] (\x,-0.3)--(\x,0.3);}
\foreach \x in {1.4,1.52,...,2.3}{\draw[rougeprepa,thick] (\x,-0.35)--(\x,0.35);}
\draw[->,rougeprepa,thick] (3.4,0) -- (4.2,0) node[midway,above]{perturbation};
\node at (2.8,-1) {Onde longitudinale dans un ressort};
\end{scope}
\end{tikzpicture}
\end{center}

\section{Célérité, durée de propagation et retard}

\subsection{Célérité d'une onde}

\begin{definitionbox}{Célérité}{celerite}
La \textbf{célérité} $v$ d'une onde est la vitesse de propagation de la perturbation dans le milieu. Si la perturbation parcourt une distance $d$ pendant une durée $\Delta t$, alors
\[
\boxed{v=\frac{d}{\Delta t}}.
\]
Dans le système international, $d$ s'exprime en mètre, $\Delta t$ en seconde et $v$ en mètre par seconde.
\end{definitionbox}

\begin{propositionbox}{Durée de propagation}{duree-propagation}
Pour une onde de célérité $v$ constante, la durée $\Delta t$ nécessaire pour parcourir une distance $d$ vaut
\[
\boxed{\Delta t=\frac{d}{v}}.
\]
\end{propositionbox}

\begin{proof}
La relation de définition est $v=d/\Delta t$. Si $v\neq 0$, on multiplie par $\Delta t$ puis on divise par $v$ :
\[
v\Delta t=d \quad \Longrightarrow \quad \Delta t=\frac{d}{v}.
\]
Cette manipulation algébrique simple doit toujours être accompagnée d'un contrôle des unités : $\si{m}/(\si{m.s^{-1}})=\si{s}$.
\end{proof}

\begin{methodbox}{Exploiter $v=d/\Delta t$}{m-vdt}
\begin{enumerate}[label=\arabic*)]
  \item Identifier la distance parcourue par la perturbation, et non la taille de l'objet qui vibre.
  \item Identifier la durée de propagation entre l'émission et la réception, ou entre deux capteurs.
  \item Convertir toutes les grandeurs dans les unités SI.
  \item Appliquer $v=d/\Delta t$, $d=v\Delta t$ ou $\Delta t=d/v$ selon l'inconnue.
  \item Conclure avec une phrase physique.
\end{enumerate}
\end{methodbox}

\begin{examplebox}{Échographie simple}{echo-simple}
Une sonde ultrasonore envoie une impulsion qui se réfléchit sur une interface située dans un tissu. L'écho revient à la sonde après $\Delta t=\SI{52}{\micro\second}$. La célérité des ultrasons dans ce tissu est $v=\SI{1540}{m.s^{-1}}$. Déterminer la profondeur de l'interface.

\textbf{Solution.} L'onde parcourt l'aller et le retour, donc une distance totale $2d$. Ainsi
\[
2d=v\Delta t \quad \Longrightarrow \quad d=\frac{v\Delta t}{2}.
\]
Avec $\Delta t=52\times 10^{-6}\,\si{s}$ :
\[
d=\frac{1540\times 52\times 10^{-6}}{2}=\SI{4.0e-2}{m}=\SI{4.0}{cm}.
\]
L'interface est située à environ $\SI{4.0}{cm}$ de la sonde.
\end{examplebox}

\begin{warningbox}{Aller simple ou aller-retour}{aller-retour}
Dans les situations d'écho, radar, sonar ou échographie, le temps mesuré est souvent un temps d'aller-retour. La distance à la cible est alors $d=v\Delta t/2$, et non $v\Delta t$.
\end{warningbox}

\subsection{Retard entre deux points}

\begin{definitionbox}{Retard}{retard}
Lorsqu'une même perturbation atteint deux points $A$ et $B$ du milieu à deux dates différentes, on appelle \textbf{retard} $\tau$ de $B$ par rapport à $A$ la durée qui sépare le passage de la perturbation en $A$ et en $B$.
Si l'onde se propage de $A$ vers $B$ à la célérité $v$, alors
\[
\boxed{\tau=\frac{AB}{v}}.
\]
\end{definitionbox}

\begin{examplebox}{Localiser un orage}{orage}
On voit l'éclair puis on entend le tonnerre $\SI{6.0}{s}$ plus tard. On néglige la durée de propagation de la lumière devant celle du son. La célérité du son dans l'air est $\SI{340}{m.s^{-1}}$.

\textbf{Solution.} La distance entre l'observateur et l'orage est
\[
d=v\Delta t=340\times 6.0=\SI{2040}{m}\approx \SI{2.0}{km}.
\]
L'orage se situe à environ $\SI{2.0}{km}$.
\end{examplebox}

\section{Ondes mécaniques périodiques}

\subsection{Périodicité temporelle}

\begin{definitionbox}{Signal périodique et période}{signal-periodique}
Un signal $s(t)$ est \textbf{périodique} de période $T$ si sa valeur se répète identiquement au bout d'une durée $T$ :
\[
\boxed{s(t+T)=s(t)} \quad \text{pour tout instant }t.
\]
La période $T$ s'exprime en seconde.
\end{definitionbox}

\begin{definitionbox}{Fréquence}{frequence}
La \textbf{fréquence} $f$ d'un signal périodique est le nombre de périodes par seconde. Elle vaut
\[
\boxed{f=\frac{1}{T}}.
\]
Elle s'exprime en hertz : $\SI{1}{Hz}=\SI{1}{s^{-1}}$.
\end{definitionbox}

\begin{center}
\begin{tikzpicture}
\begin{axis}[
width=12cm,height=5cm,
axis lines=middle,
xlabel={$t$ (s)}, ylabel={$s(t)$},
xmin=0,xmax=4.2,ymin=-1.3,ymax=1.3,
xtick={0,1,2,3,4}, ytick={-1,0,1},
grid=both]
\addplot[domain=0:4.1,samples=200,thick,bleuprepa]{sin(deg(2*pi*x))};
\draw[<->,thick,rougeprepa] (axis cs:0.25,1.12)--(axis cs:1.25,1.12) node[midway,above]{$T$};
\end{axis}
\end{tikzpicture}
\end{center}

\subsection{Périodicité spatiale et longueur d'onde}

\begin{definitionbox}{Onde mécanique périodique}{onde-periodique}
Une \textbf{onde mécanique périodique} est la propagation d'une perturbation mécanique périodique produite par une source périodique.
\end{definitionbox}

\begin{definitionbox}{Longueur d'onde}{lambda}
La \textbf{longueur d'onde} $\lambda$ est la plus petite distance séparant deux points du milieu qui vibrent dans le même état au même instant. Elle correspond aussi à la distance parcourue par l'onde pendant une période $T$.
\end{definitionbox}

\begin{retenirbox}{Double périodicité}{double-periodicite}
Une onde progressive périodique présente deux périodicités :
\begin{itemize}
  \item une \textbf{périodicité temporelle} : en un point donné, le signal se répète toutes les durées $T$ ;
  \item une \textbf{périodicité spatiale} : à un instant donné, l'aspect du milieu se répète toutes les distances $\lambda$.
\end{itemize}
\end{retenirbox}

\begin{warningbox}{Ne pas confondre $T$ et $\lambda$}{confusion-T-lambda}
La période $T$ est une durée en seconde. La longueur d'onde $\lambda$ est une longueur en mètre. On ne lit donc pas $T$ et $\lambda$ sur le même type de graphique : $T$ sur un graphique temporel, $\lambda$ sur un profil spatial.
\end{warningbox}

\subsection{Relation entre célérité, période, fréquence et longueur d'onde}

\begin{theorembox}{Relation fondamentale des ondes périodiques}{relation-fondamentale}
Pour une onde périodique de célérité $v$, de période $T$, de fréquence $f$ et de longueur d'onde $\lambda$, on a
\[
\boxed{v=\frac{\lambda}{T}} \qquad \text{et donc} \qquad \boxed{v=\lambda f}.
\]
\end{theorembox}

\begin{proof}
Par définition, pendant une période $T$, la perturbation reproduit le même état vibratoire. Le motif spatial de l'onde se décale alors d'une longueur d'onde $\lambda$. La distance parcourue pendant la durée $T$ est donc $\lambda$. En utilisant la définition de la célérité,
\[
v=\frac{\text{distance parcourue}}{\text{durée de parcours}}=\frac{\lambda}{T}.
\]
Comme $f=1/T$, on obtient immédiatement $v=\lambda f$.
\end{proof}

\begin{examplebox}{Onde sur une corde}{corde-lambda}
Une corde est excitée périodiquement à la fréquence $f=\SI{20}{Hz}$. Sur une photographie du profil de la corde, on mesure $\lambda=\SI{0.40}{m}$. La célérité est donc
\[
v=\lambda f=0.40\times 20=\SI{8.0}{m.s^{-1}}.
\]
L'onde se propage le long de la corde à $\SI{8.0}{m.s^{-1}}$.
\end{examplebox}

\begin{methodbox}{Déterminer $T$, $\lambda$ et $v$ à partir de graphiques}{graphiques}
\begin{enumerate}[label=\arabic*)]
  \item Sur un graphique $s(t)$ à un point fixé, mesurer la durée entre deux maxima successifs : c'est $T$.
  \item Sur un graphique $s(x)$ à un instant fixé, mesurer la distance entre deux maxima successifs : c'est $\lambda$.
  \item En déduire $f=1/T$.
  \item Calculer $v=\lambda/T=\lambda f$.
\end{enumerate}
\end{methodbox}

\section{Ondes sinusoïdales et modélisation mathématique}

\subsection{Signal sinusoïdal}

\begin{definitionbox}{Signal sinusoïdal}{sinusoidal}
Un signal sinusoïdal est un signal qui peut s'écrire, à une phase près,
\[
s(t)=A\cos\left(\frac{2\pi}{T}t+\varphi\right),
\]
où $A$ est l'amplitude, $T$ la période et $\varphi$ la phase initiale.
\end{definitionbox}

\begin{retenirbox}{Amplitude, période, phase}{amp-phase}
\begin{itemize}
  \item L'amplitude $A$ donne l'écart maximal à la valeur moyenne.
  \item La période $T$ donne la durée d'un motif complet.
  \item La phase $\varphi$ décale le signal horizontalement, sans changer son amplitude ni sa période.
\end{itemize}
\end{retenirbox}

\subsection{Expression simplifiée d'une onde progressive sinusoïdale}

À un niveau de première, on peut présenter une onde sinusoïdale se propageant vers les $x$ croissants sous la forme
\[
y(x,t)=A\cos\left(2\pi\left(\frac{t}{T}-\frac{x}{\lambda}\right)+\varphi\right).
\]
Cette expression n'est pas exigée comme formule à apprendre, mais elle éclaire la double périodicité.

\begin{propositionbox}{Double périodicité à partir du cosinus}{preuve-double}
La fonction
\[
y(x,t)=A\cos\left(2\pi\left(\frac{t}{T}-\frac{x}{\lambda}\right)\right)
\]
est périodique de période temporelle $T$ et de période spatiale $\lambda$.
\end{propositionbox}

\begin{proof}
À position $x$ fixée,
\[
y(x,t+T)=A\cos\left(2\pi\left(\frac{t+T}{T}-\frac{x}{\lambda}\right)\right)
=A\cos\left(2\pi\left(\frac{t}{T}-\frac{x}{\lambda}+1\right)\right).
\]
Or $\cos(\theta+2\pi)=\cos\theta$, donc $y(x,t+T)=y(x,t)$.

À date $t$ fixée,
\[
y(x+\lambda,t)=A\cos\left(2\pi\left(\frac{t}{T}-\frac{x+\lambda}{\lambda}\right)\right)
=A\cos\left(2\pi\left(\frac{t}{T}-\frac{x}{\lambda}-1\right)\right).
\]
Comme $\cos(\theta-2\pi)=\cos\theta$, on a $y(x+\lambda,t)=y(x,t)$.
\end{proof}

\subsection{Capacité numérique : simuler une onde périodique}

\begin{methodbox}{Simulation numérique d'une onde progressive}{simulation-onde}
Pour simuler une onde progressive sinusoïdale, on peut :
\begin{enumerate}[label=\arabic*)]
  \item choisir $A$, $T$, $\lambda$ et $v=\lambda/T$ ;
  \item créer une liste de positions $x$ ;
  \item calculer $y(x,t)$ à différentes dates $t$ ;
  \item représenter les profils successifs pour visualiser la propagation.
\end{enumerate}
Un exemple de formule utilisable en Python est :
\[
y=A\cos\left(2\pi\left(\frac{t}{T}-\frac{x}{\lambda}\right)\right).
\]
\end{methodbox}

\section{Exercices progressifs — Ondes mécaniques}

\begin{exercice}[Application directe]
Une onde se propage sur une corde. Une perturbation met $\SI{0.80}{s}$ pour parcourir $\SI{2.4}{m}$.
\begin{enumerate}[label=\alph*)]
  \item Calculer la célérité de l'onde.
  \item Combien de temps met-elle pour parcourir $\SI{6.0}{m}$ dans les mêmes conditions ?
\end{enumerate}
\end{exercice}

\begin{exercice}[Retard entre deux capteurs]
Deux microphones $M_1$ et $M_2$ sont séparés de $\SI{1.70}{m}$. Un claquement est produit près de $M_1$. Le signal est reçu par $M_2$ avec un retard $\tau=\SI{5.0}{ms}$.
\begin{enumerate}[label=\alph*)]
  \item Calculer la célérité du son mesurée.
  \item Comparer à la valeur usuelle $\SI{340}{m.s^{-1}}$.
\end{enumerate}
\end{exercice}

\begin{exercice}[Période, fréquence, longueur d'onde]
Une onde périodique se propage à la surface de l'eau. La période de la source est $T=\SI{0.25}{s}$ et la distance entre deux crêtes successives est $\lambda=\SI{12}{cm}$.
\begin{enumerate}[label=\alph*)]
  \item Calculer la fréquence de la source.
  \item Calculer la célérité de l'onde.
\end{enumerate}
\end{exercice}

\begin{exercice}[Graphique temporel]
En un point $M$, un capteur mesure un signal périodique. On observe 5 périodes complètes en $\SI{0.040}{s}$.
\begin{enumerate}[label=\alph*)]
  \item Déterminer la période $T$.
  \item Déterminer la fréquence $f$.
  \item Si $\lambda=\SI{8.5}{cm}$, calculer la célérité.
\end{enumerate}
\end{exercice}

\begin{exercice}[Exercice bilan — échographie]
Une sonde émet une impulsion ultrasonore dans un matériau. Deux échos sont reçus : le premier après $\SI{28}{\micro\second}$, le second après $\SI{64}{\micro\second}$. La célérité des ultrasons dans le matériau est $\SI{2500}{m.s^{-1}}$.
\begin{enumerate}[label=\alph*)]
  \item Expliquer pourquoi il faut diviser par deux la distance calculée à partir du temps d'écho.
  \item Déterminer les profondeurs des deux interfaces réfléchissantes.
  \item Déterminer l'épaisseur de la couche comprise entre les deux interfaces.
\end{enumerate}
\end{exercice}

\newpage

%============================================================
\part{La lumière : images et couleurs}
%============================================================

\section{Lentilles minces convergentes et formation des images}

\subsection{Modèle de la lentille mince convergente}

\begin{definitionbox}{Lentille mince convergente}{lentille}
Une \textbf{lentille mince convergente} est un système optique qui fait converger un faisceau de rayons parallèles à son axe optique vers un point appelé \textbf{foyer image} $F'$. Sa distance focale image est notée $f'=OF'$.
\end{definitionbox}

\begin{retenirbox}{Conventions algébriques}{conv-optique}
On utilise un axe optique orienté dans le sens de propagation de la lumière. Les distances sont algébriques :
\begin{itemize}
  \item $\overline{OA}$ est négatif si l'objet $A$ est avant la lentille ;
  \item $\overline{OA'}$ est positif si l'image $A'$ est après la lentille ;
  \item pour une lentille convergente, $f'=\overline{OF'}>0$.
\end{itemize}
\end{retenirbox}

\begin{center}
\begin{tikzpicture}[scale=1.1]
\draw[->] (-4,0)--(5,0) node[right]{axe optique};
\draw[thick,bleuprepa] (0,-1.6)--(0,1.6) node[above]{lentille};
\draw[<->,bleuprepa] (-0.25,-1.3)--(0.25,-1.3);
\node[below] at (0,0){$O$};
\fill (-1.5,0) circle(1.5pt) node[below]{$F$};
\fill (1.5,0) circle(1.5pt) node[below]{$F'$};
\draw[very thick,->,rougeprepa] (-3,0)--(-3,1.2) node[above]{$B$};
\node[below] at (-3,0){$A$};
\draw[thick,->] (-3,1.2)--(0,1.2)--(3.2,-1.36);
\draw[thick,->] (-3,1.2)--(0,0)--(3.2,-1.28);
\draw[very thick,->,vertprepa] (3.2,0)--(3.2,-1.28) node[below]{$B'$};
\node[above] at (3.2,0){$A'$};
\end{tikzpicture}
\end{center}

\subsection{Relation de conjugaison}

\begin{theorembox}{Relation de conjugaison de Descartes}{descartes}
Pour une lentille mince convergente de centre optique $O$ et de distance focale image $f'$, un objet ponctuel $A$ situé sur l'axe optique et son image $A'$ vérifient
\[
\boxed{\frac{1}{\overline{OA'}}-\frac{1}{\overline{OA}}=\frac{1}{f'}}.
\]
\end{theorembox}

\begin{remarquebox}{Statut de la relation}{statut-descartes}
En première, cette relation est fournie et utilisée. On peut cependant en donner une justification géométrique par les rayons remarquables et le théorème de Thalès, ce qui aide à comprendre le rôle des grandeurs algébriques.
\end{remarquebox}

\begin{methodbox}{Trouver la position de l'image}{position-image}
\begin{enumerate}[label=\arabic*)]
  \item Choisir l'orientation de l'axe optique et écrire les grandeurs algébriques.
  \item Identifier $\overline{OA}$ et $f'$.
  \item Utiliser
  \[
  \frac{1}{\overline{OA'}}=\frac{1}{f'}+\frac{1}{\overline{OA}}.
  \]
  \item Inverser pour obtenir $\overline{OA'}$.
  \item Interpréter le signe : image réelle si $\overline{OA'}>0$, virtuelle si $\overline{OA'}<0$.
\end{enumerate}
\end{methodbox}

\subsection{Grandissement}

\begin{definitionbox}{Grandissement}{grandissement}
Le \textbf{grandissement} $\gamma$ d'une lentille est défini par
\[
\boxed{\gamma=\frac{\overline{A'B'}}{\overline{AB}}=\frac{\overline{OA'}}{\overline{OA}}}.
\]
Il compare la taille algébrique de l'image à celle de l'objet.
\end{definitionbox}

\begin{propositionbox}{Interprétation de $\gamma$}{interpret-gamma}
\begin{itemize}
  \item Si $|\gamma|>1$, l'image est agrandie.
  \item Si $|\gamma|<1$, l'image est réduite.
  \item Si $\gamma<0$, l'image est renversée par rapport à l'objet.
  \item Si $\gamma>0$, l'image est droite.
\end{itemize}
\end{propositionbox}

\begin{examplebox}{Image réelle par une lentille convergente}{image-reelle}
Un objet $AB$ de hauteur $\SI{2.0}{cm}$ est placé à $\SI{30}{cm}$ devant une lentille convergente de distance focale $f'=\SI{10}{cm}$.

\textbf{Données algébriques.} $\overline{OA}=\SI{-30}{cm}$ et $f'=\SI{10}{cm}$.

\textbf{Position de l'image.}
\[
\frac{1}{\overline{OA'}}=\frac{1}{10}+\frac{1}{-30}=\frac{3-1}{30}=\frac{2}{30}=\frac{1}{15}.
\]
Donc $\overline{OA'}=\SI{15}{cm}$. L'image est réelle, située après la lentille.

\textbf{Grandissement.}
\[
\gamma=\frac{\overline{OA'}}{\overline{OA}}=\frac{15}{-30}=-0.50.
\]
L'image est renversée et deux fois plus petite. Sa taille est $A'B'=\gamma AB=-\SI{1.0}{cm}$, donc sa hauteur géométrique vaut $\SI{1.0}{cm}$.
\end{examplebox}

\begin{warningbox}{Oublier les signes}{signes-optique}
L'erreur la plus fréquente est d'utiliser $OA=+\SI{30}{cm}$ pour un objet réel situé avant la lentille. Avec la convention usuelle, il faut écrire $\overline{OA}=-\SI{30}{cm}$.
\end{warningbox}

\subsection{Mise au point et distance focale}

\begin{definitionbox}{Mise au point}{mise-au-point}
Réaliser une \textbf{mise au point}, c'est ajuster le système optique pour que l'image de l'objet observé se forme sur le capteur ou l'écran.
\end{definitionbox}

Pour faire une mise au point, on peut modifier :
\begin{itemize}
  \item la distance entre l'objet, la lentille et l'écran ;
  \item la distance focale de la lentille, comme dans certains dispositifs à focale variable.
\end{itemize}

\begin{methodbox}{Estimer une distance focale}{estimer-focale}
Pour estimer expérimentalement la distance focale d'une lentille convergente :
\begin{enumerate}[label=\arabic*)]
  \item viser un objet très éloigné, comme une fenêtre ou un bâtiment lointain ;
  \item déplacer un écran jusqu'à obtenir une image nette ;
  \item mesurer la distance entre la lentille et l'écran.
\end{enumerate}
Pour un objet très lointain, les rayons incidents sont presque parallèles : l'image se forme approximativement dans le plan focal image, donc $OA'\approx f'$.
\end{methodbox}

\section{Couleurs, synthèses et couleur des objets}

\subsection{Lumière blanche et couleurs complémentaires}

\begin{definitionbox}{Lumière blanche}{blanche}
Une lumière blanche est une lumière contenant un ensemble de radiations visibles. Elle peut être décomposée par un prisme ou un réseau en un spectre de couleurs.
\end{definitionbox}

\begin{definitionbox}{Couleurs complémentaires}{complementaires}
Deux couleurs lumineuses sont dites \textbf{complémentaires} si leur superposition additive peut donner une sensation de blanc.
\end{definitionbox}

\begin{center}
\begin{tikzpicture}[scale=0.9]
\draw[fill=red!55,opacity=.75] (0,0) circle (1.4);
\draw[fill=green!55,opacity=.75] (1.4,0) circle (1.4);
\draw[fill=blue!55,opacity=.75] (0.7,1.2) circle (1.4);
\node at (-0.8,-0.5){R};
\node at (2.2,-0.5){V};
\node at (0.7,2.1){B};
\node at (0.7,0.45){blanc};
\node at (0.7,-1.7){Synthèse additive : lumières};
\end{tikzpicture}
\end{center}

\subsection{Synthèse additive et synthèse soustractive}

\begin{definitionbox}{Synthèse additive}{additive}
La \textbf{synthèse additive} consiste à superposer des lumières colorées. Les couleurs primaires de ce modèle sont généralement le rouge, le vert et le bleu.
\end{definitionbox}

\begin{definitionbox}{Synthèse soustractive}{soustractive}
La \textbf{synthèse soustractive} concerne des filtres, pigments ou encres qui absorbent certaines radiations de la lumière incidente et en transmettent ou diffusent d'autres. Les couleurs primaires de ce modèle sont généralement le cyan, le magenta et le jaune.
\end{definitionbox}

\begin{retenirbox}{Choisir le bon modèle}{modele-couleurs}
\begin{itemize}
  \item Écran, projecteurs, pixels lumineux : synthèse additive.
  \item Peinture, encre, filtres, objets éclairés : synthèse soustractive et phénomènes d'absorption, diffusion, transmission.
\end{itemize}
\end{retenirbox}

\begin{tabularx}{\textwidth}{>{\bfseries}lX}
Absorption & une partie de la lumière incidente est absorbée par la matière.\\
Diffusion & une partie de la lumière incidente est renvoyée dans plusieurs directions.\\
Transmission & une partie de la lumière traverse le milieu ou le filtre.\\
\end{tabularx}

\begin{examplebox}{Objet rouge éclairé en lumière blanche}{objet-rouge}
Un objet rouge éclairé en lumière blanche apparaît rouge car il diffuse principalement les radiations rouges et absorbe une grande partie des autres radiations visibles.

S'il est éclairé par une lumière bleue seule, il peut apparaître sombre, car il ne reçoit presque pas de radiation rouge à diffuser.
\end{examplebox}

\begin{warningbox}{La couleur n'est pas une propriété absolue}{couleur-non-absolue}
Un objet n'a pas une couleur indépendante de l'éclairage. La couleur perçue dépend de la lumière incidente, des propriétés d'absorption/diffusion/transmission de l'objet, et de l'observateur.
\end{warningbox}

\subsection{Vision des couleurs et trichromie}

\begin{definitionbox}{Trichromie}{trichromie}
La vision humaine des couleurs repose sur trois types de cônes sensibles à des domaines différents du visible. Le cerveau interprète les signaux issus de ces cônes pour produire une sensation colorée.
\end{definitionbox}

\begin{remarquebox}{Lien avec les écrans}{ecrans}
Un écran n'a pas besoin d'émettre toutes les radiations visibles pour donner une impression de couleur. Il règle l'intensité de sous-pixels rouges, verts et bleus afin de stimuler les cônes de l'œil comme le ferait une lumière d'une certaine couleur.
\end{remarquebox}

\section{Exercices progressifs — Images et couleurs}

\begin{exercice}[Lentille convergente]
Un objet de hauteur $\SI{3.0}{cm}$ est placé à $\SI{40}{cm}$ devant une lentille convergente de distance focale $\SI{20}{cm}$.
\begin{enumerate}[label=\alph*)]
  \item Écrire $\overline{OA}$ et $f'$ avec les signes corrects.
  \item Déterminer $\overline{OA'}$.
  \item Calculer le grandissement.
  \item Donner les caractéristiques de l'image.
\end{enumerate}
\end{exercice}

\begin{exercice}[Image virtuelle]
Un objet est placé à $\SI{6.0}{cm}$ devant une lentille convergente de distance focale $\SI{10}{cm}$.
\begin{enumerate}[label=\alph*)]
  \item Calculer la position de l'image.
  \item L'image est-elle réelle ou virtuelle ?
  \item Calculer le grandissement et interpréter.
\end{enumerate}
\end{exercice}

\begin{exercice}[Filtres colorés]
Une lumière blanche traverse successivement un filtre jaune puis un filtre cyan.
\begin{enumerate}[label=\alph*)]
  \item Rappeler quelles couleurs sont transmises par un filtre jaune dans un modèle simple RVB.
  \item Rappeler quelles couleurs sont transmises par un filtre cyan.
  \item Prévoir la couleur finale transmise.
\end{enumerate}
\end{exercice}

\begin{exercice}[Couleur d'un objet]
Une pomme rouge est éclairée successivement par une lumière blanche, une lumière rouge, puis une lumière verte.
\begin{enumerate}[label=\alph*)]
  \item Quelle couleur perçoit-on sous lumière blanche ?
  \item Quelle couleur perçoit-on sous lumière rouge ?
  \item Quelle couleur perçoit-on sous lumière verte ? Justifier.
\end{enumerate}
\end{exercice}

\newpage

%============================================================
\part{Modèles ondulatoire et particulaire de la lumière}
%============================================================

\section{La lumière comme onde électromagnétique}

\subsection{Domaines spectraux}

\begin{definitionbox}{Onde électromagnétique}{onde-em}
Une \textbf{onde électromagnétique} est une onde pouvant se propager dans le vide. La lumière visible est une petite partie du spectre des ondes électromagnétiques.
\end{definitionbox}

\begin{retenirbox}{Célérité de la lumière dans le vide}{c}
Dans le vide, la lumière se propage à la célérité
\[
\boxed{c=\SI{3.00e8}{m.s^{-1}}}.
\]
Dans l'air, on utilise souvent la même valeur en première.
\end{retenirbox}

\begin{center}
\begin{tikzpicture}[xscale=1.15,yscale=1]
\draw[->,thick] (0,0)--(12,0) node[right]{fréquence croissante};
\foreach \x/\lab in {0.5/radio,2.5/micro-ondes,4.5/infrarouge,6/visible,7.5/ultraviolet,9.5/rayons X,11/gamma}{
  \draw (\x,0.12)--(\x,-0.12);
  \node[rotate=35,anchor=west] at (\x,-0.35){\lab};
}
\draw[fill=violet!35] (5.55,0.25) rectangle (6.45,0.75);
\node at (6,1.05){visible};
\end{tikzpicture}
\end{center}

\begin{table}[h]
\centering
\begin{tabularx}{\textwidth}{>{\bfseries}lXX}
\toprule
Domaine & Ordre de grandeur de longueur d'onde & Applications typiques\\
\midrule
Ondes radio & m à km & radiodiffusion, communications\\
Micro-ondes & mm à cm & Wi-Fi, four micro-ondes, radar\\
Infrarouge & $\mu$m & télécommandes, imagerie thermique\\
Visible & $\SI{400}{nm}$ à $\SI{800}{nm}$ environ & vision, photographie\\
Ultraviolet & dizaines à centaines de nm & fluorescence, stérilisation\\
Rayons X & autour du nm et moins & radiographie médicale\\
Gamma & très petites longueurs d'onde & phénomènes nucléaires, médecine nucléaire\\
\bottomrule
\end{tabularx}
\end{table}

\subsection{Relation entre longueur d'onde, fréquence et célérité}

\begin{theorembox}{Relation fondamentale pour la lumière}{lambda-c-nu}
Dans le vide ou dans l'air, une onde lumineuse de fréquence $\nu$ et de longueur d'onde $\lambda$ vérifie
\[
\boxed{c=\lambda\nu} \qquad \text{donc} \qquad \boxed{\lambda=\frac{c}{\nu}}.
\]
\end{theorembox}

\begin{examplebox}{Fréquence d'une radiation rouge}{rouge-freq}
Une radiation rouge a pour longueur d'onde dans l'air $\lambda=\SI{650}{nm}$.
\[
\nu=\frac{c}{\lambda}=\frac{3.00\times 10^8}{650\times 10^{-9}}=\SI{4.62e14}{Hz}.
\]
La fréquence est de l'ordre de $\SI{4.6e14}{Hz}$.
\end{examplebox}

\begin{warningbox}{Fréquence et longueur d'onde}{freq-lambda}
Quand la fréquence augmente, la longueur d'onde diminue puisque $\lambda=c/\nu$. Les ultraviolets ont donc une fréquence plus grande et une longueur d'onde plus petite que le visible rouge.
\end{warningbox}

\section{Le photon et l'énergie lumineuse}

\subsection{Modèle particulaire}

\begin{definitionbox}{Photon}{photon}
Dans le modèle particulaire, la lumière peut être décrite comme un flux de particules appelées \textbf{photons}. Chaque photon transporte une énergie qui dépend de la fréquence de la radiation.
\end{definitionbox}

\begin{theorembox}{Énergie d'un photon}{energie-photon}
L'énergie $E$ d'un photon de fréquence $\nu$ vaut
\[
\boxed{E=h\nu}
\]
où $h$ est la constante de Planck :
\[
h=\SI{6.63e-34}{J.s}.
\]
En utilisant $c=\lambda\nu$, on obtient aussi
\[
\boxed{E=\frac{hc}{\lambda}}.
\]
\end{theorembox}

\begin{proof}
La relation $E=h\nu$ est une loi du modèle quantique de la lumière. À partir de la relation ondulatoire $c=\lambda\nu$, on a $\nu=c/\lambda$. En remplaçant dans $E=h\nu$, on obtient
\[
E=h\frac{c}{\lambda}=\frac{hc}{\lambda}.
\]
Cette démonstration montre la cohérence entre le modèle ondulatoire et le modèle particulaire : ils ne s'excluent pas, ils éclairent deux aspects différents de la lumière.
\end{proof}

\begin{examplebox}{Énergie d'un photon vert}{photon-vert}
Un photon vert a une longueur d'onde $\lambda=\SI{530}{nm}$.
\[
E=\frac{hc}{\lambda}=\frac{6.63\times 10^{-34}\times 3.00\times 10^8}{530\times 10^{-9}}=\SI{3.75e-19}{J}.
\]
En électronvolt, avec $\SI{1}{eV}=\SI{1.60e-19}{J}$,
\[
E=\frac{3.75\times 10^{-19}}{1.60\times 10^{-19}}=\SI{2.34}{eV}.
\]
\end{examplebox}

\begin{retenirbox}{Ordres de grandeur}{ordre-photon}
Les photons visibles ont des énergies de l'ordre de quelques électronvolts. Les photons ultraviolets sont plus énergétiques que les photons visibles, eux-mêmes plus énergétiques que les photons infrarouges.
\end{retenirbox}

\section{Interaction lumière-matière et niveaux d'énergie}

\subsection{Quantification des niveaux d'énergie}

\begin{definitionbox}{Niveaux d'énergie quantifiés}{niveaux}
Dans un atome, l'énergie ne peut prendre que certaines valeurs discrètes appelées \textbf{niveaux d'énergie}. On dit que l'énergie est quantifiée.
\end{definitionbox}

\begin{propositionbox}{Absorption et émission}{abs-em}
Un atome peut :
\begin{itemize}
  \item \textbf{absorber} un photon si l'énergie du photon correspond exactement à l'écart entre deux niveaux d'énergie ;
  \item \textbf{émettre} un photon lorsqu'il passe d'un niveau d'énergie plus élevé à un niveau plus bas.
\end{itemize}
Dans les deux cas,
\[
\boxed{|\Delta E|=h\nu=\frac{hc}{\lambda}}.
\]
\end{propositionbox}

\begin{center}
\begin{tikzpicture}[scale=1]
\draw[thick] (0,0)--(4,0) node[right]{$E_1$};
\draw[thick] (0,1.4)--(4,1.4) node[right]{$E_2$};
\draw[thick] (0,3.0)--(4,3.0) node[right]{$E_3$};
\draw[->,rougeprepa,thick] (1,0.05)--(1,1.35) node[midway,left]{absorption};
\draw[->,bleuprepa,thick] (2.7,3.0)--(2.7,1.45) node[midway,right]{émission};
\node at (2,-0.7){Diagramme qualitatif de niveaux d'énergie};
\end{tikzpicture}
\end{center}

\begin{methodbox}{Exploiter un diagramme de niveaux d'énergie}{diagramme-niveaux}
\begin{enumerate}[label=\arabic*)]
  \item Identifier le niveau initial et le niveau final.
  \item Calculer l'écart énergétique $|\Delta E|=|E_f-E_i|$.
  \item Convertir l'énergie en joules si nécessaire.
  \item Utiliser $|\Delta E|=h\nu$ ou $|\Delta E|=hc/\lambda$.
  \item Conclure : absorption si l'atome gagne de l'énergie, émission s'il en perd.
\end{enumerate}
\end{methodbox}

\begin{examplebox}{Transition atomique}{transition}
Un atome passe d'un niveau excité $E_2=-\SI{3.40}{eV}$ à un niveau plus bas $E_1=-\SI{13.6}{eV}$.

\textbf{Écart d'énergie.}
\[
\Delta E=E_2-E_1=(-3.40)-(-13.6)=\SI{10.2}{eV}.
\]
En joules :
\[
\Delta E=10.2\times 1.60\times 10^{-19}=\SI{1.63e-18}{J}.
\]

\textbf{Longueur d'onde émise.}
\[
\lambda=\frac{hc}{\Delta E}=\frac{6.63\times 10^{-34}\times 3.00\times 10^8}{1.63\times 10^{-18}}=\SI{1.22e-7}{m}=\SI{122}{nm}.
\]
Cette radiation appartient à l'ultraviolet.
\end{examplebox}

\subsection{Spectres d'émission et d'absorption}

\begin{definitionbox}{Spectre}{spectre}
Un \textbf{spectre} représente la répartition d'une lumière selon ses longueurs d'onde ou ses fréquences. Il peut être continu ou composé de raies.
\end{definitionbox}

\begin{propositionbox}{Signature d'une entité chimique}{signature}
Les raies d'émission ou d'absorption d'une entité chimique dépendent des écarts entre ses niveaux d'énergie. Elles constituent donc une signature permettant d'identifier cette entité.
\end{propositionbox}

\begin{warningbox}{Spectre continu et spectre de raies}{spectres-confusion}
Un spectre continu contient toutes les longueurs d'onde d'un intervalle. Un spectre de raies ne contient que certaines longueurs d'onde. Les raies traduisent la quantification des niveaux d'énergie.
\end{warningbox}

\section{Exercices progressifs — Modèles de la lumière}

\begin{exercice}[Domaine spectral]
Une onde électromagnétique a une fréquence $\nu=\SI{2.4e9}{Hz}$.
\begin{enumerate}[label=\alph*)]
  \item Calculer sa longueur d'onde dans l'air.
  \item Identifier le domaine spectral correspondant.
  \item Citer une application courante de ce domaine.
\end{enumerate}
\end{exercice}

\begin{exercice}[Photon visible]
Calculer l'énergie en joules puis en électronvolts d'un photon de longueur d'onde $\lambda=\SI{450}{nm}$. On donne $h=\SI{6.63e-34}{J.s}$, $c=\SI{3.00e8}{m.s^{-1}}$ et $\SI{1}{eV}=\SI{1.60e-19}{J}$.
\end{exercice}

\begin{exercice}[Diagramme de niveaux]
Un atome possède trois niveaux d'énergie :
\[
E_1=-\SI{5.0}{eV},\qquad E_2=-\SI{3.0}{eV},\qquad E_3=-\SI{1.0}{eV}.
\]
\begin{enumerate}[label=\alph*)]
  \item Calculer les énergies des photons pouvant être émis.
  \item Déterminer la longueur d'onde correspondant à la transition $E_3\to E_1$.
  \item Cette radiation est-elle visible ?
\end{enumerate}
\end{exercice}

\newpage

%============================================================
\part{Corrections détaillées des exercices}
%============================================================

\setcounter{correction}{0}

\begin{correction}[Application directe]
\begin{enumerate}[label=\alph*)]
  \item La célérité vaut
  \[
  v=\frac{d}{\Delta t}=\frac{2.4}{0.80}=\SI{3.0}{m.s^{-1}}.
  \]
  \item Dans les mêmes conditions, $v$ est constante, donc
  \[
  \Delta t=\frac{d}{v}=\frac{6.0}{3.0}=\SI{2.0}{s}.
  \]
\end{enumerate}
\end{correction}

\begin{correction}[Retard entre deux capteurs]
\begin{enumerate}[label=\alph*)]
  \item Le retard est $\tau=\SI{5.0}{ms}=\SI{5.0e-3}{s}$. Donc
  \[
  v=\frac{M_1M_2}{\tau}=\frac{1.70}{5.0\times 10^{-3}}=\SI{340}{m.s^{-1}}.
  \]
  \item La valeur est en accord avec la célérité usuelle du son dans l'air.
\end{enumerate}
\end{correction}

\begin{correction}[Période, fréquence, longueur d'onde]
\begin{enumerate}[label=\alph*)]
  \item $f=1/T=1/0.25=\SI{4.0}{Hz}$.
  \item $\lambda=\SI{12}{cm}=\SI{0.12}{m}$, donc
  \[
  v=\lambda f=0.12\times 4.0=\SI{0.48}{m.s^{-1}}.
  \]
\end{enumerate}
\end{correction}

\begin{correction}[Graphique temporel]
\begin{enumerate}[label=\alph*)]
  \item Si 5 périodes durent $\SI{0.040}{s}$, alors
  \[
  T=\frac{0.040}{5}=\SI{8.0e-3}{s}.
  \]
  \item $f=1/T=1/(8.0\times 10^{-3})=\SI{125}{Hz}$.
  \item $\lambda=\SI{8.5}{cm}=\SI{0.085}{m}$, donc
  \[
  v=\lambda f=0.085\times 125=\SI{10.6}{m.s^{-1}}.
  \]
\end{enumerate}
\end{correction}

\begin{correction}[Exercice bilan — échographie]
\begin{enumerate}[label=\alph*)]
  \item Le temps d'écho correspond au trajet aller jusqu'à l'interface puis retour jusqu'à la sonde. La distance totale parcourue est donc $2d$.
  \item Pour le premier écho :
  \[
  d_1=\frac{v\Delta t_1}{2}=\frac{2500\times 28\times 10^{-6}}{2}=\SI{3.5e-2}{m}=\SI{3.5}{cm}.
  \]
  Pour le second :
  \[
  d_2=\frac{2500\times 64\times 10^{-6}}{2}=\SI{8.0e-2}{m}=\SI{8.0}{cm}.
  \]
  \item L'épaisseur est $d_2-d_1=\SI{4.5}{cm}$.
\end{enumerate}
\end{correction}

\begin{correction}[Lentille convergente]
\begin{enumerate}[label=\alph*)]
  \item L'objet est devant la lentille : $\overline{OA}=\SI{-40}{cm}$. La lentille est convergente : $f'=\SI{20}{cm}$.
  \item
  \[
  \frac{1}{\overline{OA'}}=\frac{1}{20}+\frac{1}{-40}=\frac{2-1}{40}=\frac{1}{40}.
  \]
  Donc $\overline{OA'}=\SI{40}{cm}$.
  \item
  \[
  \gamma=\frac{\overline{OA'}}{\overline{OA}}=\frac{40}{-40}=-1.
  \]
  \item L'image est réelle, située après la lentille, renversée et de même taille que l'objet.
\end{enumerate}
\end{correction}

\begin{correction}[Image virtuelle]
\begin{enumerate}[label=\alph*)]
  \item $\overline{OA}=\SI{-6.0}{cm}$ et $f'=\SI{10}{cm}$. Alors
  \[
  \frac{1}{\overline{OA'}}=\frac{1}{10}+\frac{1}{-6.0}=\frac{3-5}{30}=-\frac{2}{30}=-\frac{1}{15}.
  \]
  Donc $\overline{OA'}=\SI{-15}{cm}$.
  \item L'image est virtuelle car elle se situe du même côté que l'objet.
  \item
  \[
  \gamma=\frac{-15}{-6.0}=2.5.
  \]
  L'image est droite et agrandie.
\end{enumerate}
\end{correction}

\begin{correction}[Filtres colorés]
Dans un modèle simple RVB :
\begin{enumerate}[label=\alph*)]
  \item Un filtre jaune transmet le rouge et le vert, et absorbe le bleu.
  \item Un filtre cyan transmet le vert et le bleu, et absorbe le rouge.
  \item Après le filtre jaune, il reste rouge + vert. Le filtre cyan absorbe le rouge et transmet le vert. La lumière finale est donc verte.
\end{enumerate}
\end{correction}

\begin{correction}[Couleur d'un objet]
\begin{enumerate}[label=\alph*)]
  \item Sous lumière blanche, la pomme diffuse surtout le rouge : elle apparaît rouge.
  \item Sous lumière rouge, elle reçoit du rouge et le diffuse : elle apparaît rouge, éventuellement très lumineuse.
  \item Sous lumière verte, elle reçoit peu ou pas de rouge à diffuser. Elle absorbe majoritairement le vert : elle apparaît sombre ou noire.
\end{enumerate}
\end{correction}

\begin{correction}[Domaine spectral]
\begin{enumerate}[label=\alph*)]
  \item
  \[
  \lambda=\frac{c}{\nu}=\frac{3.00\times 10^8}{2.4\times 10^9}=\SI{0.125}{m}=\SI{12.5}{cm}.
  \]
  \item Cette longueur d'onde correspond au domaine des micro-ondes.
  \item Une application courante est le Wi-Fi.
\end{enumerate}
\end{correction}

\begin{correction}[Photon visible]
\[
E=\frac{hc}{\lambda}=\frac{6.63\times 10^{-34}\times 3.00\times 10^8}{450\times 10^{-9}}=\SI{4.42e-19}{J}.
\]
En électronvolts :
\[
E=\frac{4.42\times 10^{-19}}{1.60\times 10^{-19}}=\SI{2.76}{eV}.
\]
\end{correction}

\begin{correction}[Diagramme de niveaux]
\begin{enumerate}[label=\alph*)]
  \item Les émissions possibles correspondent aux transitions descendantes :
  \[
  E_3\to E_2 : \Delta E=\SI{2.0}{eV},\qquad
  E_2\to E_1 : \Delta E=\SI{2.0}{eV},\qquad
  E_3\to E_1 : \Delta E=\SI{4.0}{eV}.
  \]
  \item Pour $E_3\to E_1$, $\Delta E=\SI{4.0}{eV}=4.0\times 1.60\times 10^{-19}=\SI{6.4e-19}{J}$.
  \[
  \lambda=\frac{hc}{\Delta E}=\frac{6.63\times 10^{-34}\times 3.00\times 10^8}{6.4\times 10^{-19}}=\SI{3.11e-7}{m}=\SI{311}{nm}.
  \]
  \item $\SI{311}{nm}$ est inférieur à $\SI{400}{nm}$ : c'est de l'ultraviolet, donc non visible pour l'œil humain.
\end{enumerate}
\end{correction}

\newpage

%============================================================
\part{Grand devoir bilan final}
%============================================================

\section*{Sujet — Devoir bilan : Ondes et signaux}
\addcontentsline{toc}{section}{Sujet — Devoir bilan final}

\begin{tcolorbox}[colframe=bleuprepa,colback=white,title={Devoir bilan final — Durée conseillée : 2 h}]
\textbf{Consignes.} Les réponses doivent être rédigées, les unités précisées et les résultats numériques donnés avec un nombre raisonnable de chiffres significatifs. Les schémas sont attendus lorsqu'ils aident au raisonnement.
\end{tcolorbox}

\subsection*{Exercice A — Onde sonore et localisation}
Deux microphones $M_1$ et $M_2$ sont distants de $\SI{2.04}{m}$. Un son bref est produit sur la droite des microphones. Il est reçu par $M_1$ puis par $M_2$ avec un retard $\tau=\SI{6.0}{ms}$.
\begin{enumerate}[label=\arabic*)]
  \item Définir une onde mécanique progressive.
  \item Calculer la célérité mesurée du son.
  \item Le son peut-il se propager dans le vide ? Justifier.
  \item On réalise ensuite une expérience d'écho contre un mur. L'écho revient après $\SI{0.18}{s}$. Déterminer la distance au mur.
\end{enumerate}

\subsection*{Exercice B — Onde périodique à la surface de l'eau}
Une source produit des rides circulaires périodiques à la surface de l'eau. On mesure $12$ intervalles entre crêtes successives sur une distance de $\SI{18.0}{cm}$. La source effectue $30$ oscillations en $\SI{10.0}{s}$.
\begin{enumerate}[label=\arabic*)]
  \item Déterminer la longueur d'onde.
  \item Déterminer la fréquence et la période.
  \item Calculer la célérité de l'onde.
  \item Expliquer la différence entre périodicité spatiale et périodicité temporelle.
\end{enumerate}

\subsection*{Exercice C — Lentille et photographie}
Un objet réel de hauteur $\SI{1.80}{m}$ est situé à $\SI{3.0}{m}$ devant une lentille mince convergente de distance focale $f'=\SI{50}{mm}$.
\begin{enumerate}[label=\arabic*)]
  \item Écrire les données algébriques en mètre.
  \item Calculer la position de l'image.
  \item Calculer le grandissement.
  \item En déduire la taille de l'image sur le capteur.
  \item L'image est-elle droite ou renversée ? réelle ou virtuelle ?
\end{enumerate}

\subsection*{Exercice D — Couleurs et filtres}
Une scène est éclairée en lumière blanche. Un tissu diffuse le rouge et le vert, mais absorbe le bleu.
\begin{enumerate}[label=\arabic*)]
  \item Quelle couleur perçoit-on ?
  \item On éclaire maintenant le tissu avec une lumière bleue seule. Quelle couleur perçoit-on ?
  \item On place devant une lumière blanche un filtre magenta puis un filtre cyan. Prévoir la couleur transmise dans un modèle RVB simple.
\end{enumerate}

\subsection*{Exercice E — Photon et niveaux d'énergie}
Une radiation lumineuse a une longueur d'onde $\lambda=\SI{486}{nm}$. On donne $h=\SI{6.63e-34}{J.s}$, $c=\SI{3.00e8}{m.s^{-1}}$, $\SI{1}{eV}=\SI{1.60e-19}{J}$.
\begin{enumerate}[label=\arabic*)]
  \item Calculer sa fréquence.
  \item Calculer l'énergie d'un photon en joules puis en électronvolts.
  \item Cette radiation peut-elle être émise lors d'une transition atomique d'écart énergétique $\SI{2.55}{eV}$ ?
  \item Expliquer pourquoi un spectre de raies permet d'identifier une entité chimique.
\end{enumerate}

\newpage

\section*{Correction détaillée du devoir bilan}
\addcontentsline{toc}{section}{Correction détaillée du devoir bilan final}

\subsection*{Correction A}
\begin{enumerate}[label=\arabic*)]
  \item Une onde mécanique progressive est la propagation d'une perturbation mécanique dans un milieu matériel, sans transport global de matière.
  \item
  \[
  v=\frac{M_1M_2}{\tau}=\frac{2.04}{6.0\times10^{-3}}=\SI{340}{m.s^{-1}}.
  \]
  \item Non. Le son est une onde mécanique : il nécessite un milieu matériel pour se propager.
  \item Le temps mesuré est un aller-retour :
  \[
  d=\frac{v\Delta t}{2}=\frac{340\times0.18}{2}=\SI{30.6}{m}\approx\SI{31}{m}.
  \]
\end{enumerate}

\subsection*{Correction B}
\begin{enumerate}[label=\arabic*)]
  \item $12$ intervalles crête-crête correspondent à $12\lambda$, donc
  \[
  \lambda=\frac{18.0}{12}=\SI{1.50}{cm}=\SI{1.50e-2}{m}.
  \]
  \item
  \[
  f=\frac{30}{10.0}=\SI{3.0}{Hz},\qquad T=\frac{1}{f}=\SI{0.333}{s}.
  \]
  \item
  \[
  v=\lambda f=1.50\times10^{-2}\times3.0=\SI{4.5e-2}{m.s^{-1}}.
  \]
  \item La périodicité temporelle se lit en un point fixe et correspond à $T$. La périodicité spatiale se lit à une date fixe et correspond à $\lambda$.
\end{enumerate}

\subsection*{Correction C}
\begin{enumerate}[label=\arabic*)]
  \item $\overline{OA}=\SI{-3.0}{m}$ et $f'=\SI{50}{mm}=\SI{5.0e-2}{m}$.
  \item
  \[
  \frac{1}{\overline{OA'}}=\frac{1}{0.050}+\frac{1}{-3.0}=20-0.333=19.667.
  \]
  Donc
  \[
  \overline{OA'}=\frac{1}{19.667}=\SI{5.08e-2}{m}=\SI{5.08}{cm}.
  \]
  \item
  \[
  \gamma=\frac{\overline{OA'}}{\overline{OA}}=\frac{0.0508}{-3.0}=-1.69\times10^{-2}.
  \]
  \item
  \[
  A'B'=\gamma AB=-1.69\times10^{-2}\times1.80=-\SI{3.05e-2}{m}.
  \]
  La taille géométrique de l'image est donc environ $\SI{3.1}{cm}$.
  \item L'image est réelle car $\overline{OA'}>0$ et renversée car $\gamma<0$.
\end{enumerate}

\subsection*{Correction D}
\begin{enumerate}[label=\arabic*)]
  \item Rouge + vert donne une sensation de jaune : le tissu apparaît jaune.
  \item Sous lumière bleue seule, le tissu absorbe le bleu et ne reçoit ni rouge ni vert à diffuser : il apparaît sombre.
  \item Un filtre magenta transmet rouge + bleu. Un filtre cyan transmet vert + bleu. La seule composante commune transmise par les deux est le bleu. La lumière finale est bleue.
\end{enumerate}

\subsection*{Correction E}
\begin{enumerate}[label=\arabic*)]
  \item
  \[
  \nu=\frac{c}{\lambda}=\frac{3.00\times10^8}{486\times10^{-9}}=\SI{6.17e14}{Hz}.
  \]
  \item
  \[
  E=h\nu=6.63\times10^{-34}\times6.17\times10^{14}=\SI{4.09e-19}{J}.
  \]
  En électronvolts :
  \[
  E=\frac{4.09\times10^{-19}}{1.60\times10^{-19}}=\SI{2.56}{eV}.
  \]
  \item Oui, à l'incertitude des arrondis près, $\SI{2.56}{eV}$ correspond à $\SI{2.55}{eV}$.
  \item Les niveaux d'énergie d'une entité sont quantifiés et propres à cette entité. Les écarts entre niveaux déterminent les longueurs d'onde des photons absorbés ou émis : le spectre de raies constitue donc une signature.
\end{enumerate}

\newpage

%============================================================
\part{Fiche de synthèse finale}
%============================================================

\section*{Fiche de synthèse — Ondes et signaux}
\addcontentsline{toc}{section}{Fiche de synthèse finale}

\begin{tcolorbox}[colframe=bleuprepa,colback=blue!3!white,title={1. Ondes mécaniques}]
\begin{itemize}[label=$\triangleright$]
  \item Une onde mécanique progressive est la propagation d'une perturbation dans un milieu matériel, sans transport global de matière.
  \item La célérité est $v=d/\Delta t$.
  \item Le retard entre deux points séparés de $AB$ vaut $\tau=AB/v$.
  \item Une onde périodique possède une période temporelle $T$ et une longueur d'onde $\lambda$.
  \item La fréquence est $f=1/T$.
  \item La relation fondamentale est $v=\lambda/T=\lambda f$.
\end{itemize}
\end{tcolorbox}

\begin{tcolorbox}[colframe=vertprepa,colback=green!4!white,title={2. Lentilles minces convergentes}]
\begin{itemize}[label=$\triangleright$]
  \item Relation de conjugaison :
  \[
  \frac{1}{\overline{OA'}}-\frac{1}{\overline{OA}}=\frac{1}{f'}.
  \]
  \item Grandissement :
  \[
  \gamma=\frac{\overline{A'B'}}{\overline{AB}}=\frac{\overline{OA'}}{\overline{OA}}.
  \]
  \item $\overline{OA'}>0$ : image réelle ; $\overline{OA'}<0$ : image virtuelle.
  \item $\gamma<0$ : image renversée ; $\gamma>0$ : image droite.
\end{itemize}
\end{tcolorbox}

\begin{tcolorbox}[colframe=orangeprepa,colback=orange!4!white,title={3. Couleurs}]
\begin{itemize}[label=$\triangleright$]
  \item Synthèse additive : superposition de lumières colorées, modèle RVB.
  \item Synthèse soustractive : filtres, pigments, encres, modèle CMJ.
  \item La couleur perçue d'un objet dépend de la lumière incidente, de l'absorption, de la diffusion, de la transmission et de l'œil.
  \item Un objet rouge n'apparaît rouge que s'il reçoit des radiations rouges qu'il peut diffuser.
\end{itemize}
\end{tcolorbox}

\begin{tcolorbox}[colframe=violetprepa,colback=violet!4!white,title={4. Lumière : modèles ondulatoire et particulaire}]
\begin{itemize}[label=$\triangleright$]
  \item Dans le vide ou l'air : $c=\lambda\nu$ avec $c=\SI{3.00e8}{m.s^{-1}}$.
  \item Énergie d'un photon : $E=h\nu=hc/\lambda$.
  \item Les niveaux d'énergie atomiques sont quantifiés.
  \item Absorption ou émission : $|\Delta E|=h\nu=hc/\lambda$.
  \item Un spectre de raies permet d'identifier une entité chimique.
\end{itemize}
\end{tcolorbox}

\begin{tcolorbox}[colframe=rougeprepa,colback=red!3!white,title={Erreurs classiques à éviter}]
\begin{enumerate}[label=\arabic*)]
  \item Confondre période $T$ et longueur d'onde $\lambda$.
  \item Oublier de convertir les millisecondes, microsecondes, nanomètres ou centimètres.
  \item Oublier le facteur $2$ dans les expériences d'écho.
  \item Utiliser des distances non algébriques dans la relation de conjugaison.
  \item Croire que la couleur d'un objet est indépendante de la lumière qui l'éclaire.
  \item Penser qu'une onde mécanique peut se propager dans le vide.
\end{enumerate}
\end{tcolorbox}

\end{document}
